%% Section 1: Introduction
%% English manuscript section.
%% -----------------------------------------------------------------------

\section{Introduction}
\label{sec:introduction}

The On-Line Encyclopedia of Integer Sequences (OEIS)~\cite{sloane1996}, as of January 2026,
catalogues 391,710 entries spanning combinatorics, number theory, algebra, and many
other branches of mathematics, making it the \emph{de facto} standard reference for integer
sequences. Each entry associates a finite integer sequence with its mathematical definition,
rendering the OEIS a uniquely machine-readable corpus of mathematical knowledge.

The long-term goal of this research is an AI system that discovers hidden number-theoretic
similarities between integer sequences arising in different mathematical fields or from
different algorithms, and thereby supports the generation of mathematical conjectures---a
goal in the same spirit as recent AI-driven discovery systems such as
AlphaProof~\cite{alphaproof2024}, the Ramanujan Machine~\cite{raayoni2021ramanujan}, and
FunSearch~\cite{romera2023funsearch} (Section~\ref{sec:background:ai-math}). A prerequisite
for this goal is a machine representation that captures the number-theoretic structure of
OEIS sequences. This paper takes a first step: we propose such a representation and measure,
through a masked-prediction task, how much arithmetic structure it actually captures.

The task we formalise is \emph{masked sequence modelling}: a random subset of positions in a
sequence is masked and the model is trained to predict each masked value from its surrounding
context. This task directly probes how well a model has internalised the arithmetic and
combinatorial laws governing integer sequences; next-term prediction is treated as a special
case and evaluated via a dedicated Solver component. The prior benchmark FACT~\cite{zurich-fact}
defines and evaluates five tasks including unmasking, establishing what a plain Transformer can
achieve, but it faces fundamental limitations in handling large integer values outside its fixed
token vocabulary and in learning multiplicative structure.

The principal difficulty of this task is the extreme heterogeneity of integer sequences.
Values range from single-digit constants to astronomically large factorials and exponential
sequences, with differences of tens of orders of magnitude within a single training corpus.
At the same time, many sequences obey \emph{residue constraints}---parity, combinatorial
patterns, or periodic remainders---that are independent of the magnitude of individual values.
The standard tokenisation approach of assigning each integer a discrete vocabulary token is
fundamentally ill-suited to this setting: it cannot represent integers outside its fixed token
vocabulary, embeds arithmetic structure in opaque token IDs, and breaks down in scale for large numbers.

We propose \textbf{IntSeqBERT}, a Transformer encoder pre-trained on the OEIS corpus via masked
sequence modelling, which overcomes the above challenges through a \emph{dual-stream} input
representation. Rather than tokenising integers, IntSeqBERT encodes each element along two
complementary axes:
\begin{itemize}
  \item \textbf{Magnitude stream}: a continuous log\-scale embedding of the
        absolute value, capturing growth behaviour and scale.
  \item \textbf{Modulo stream}: sin/cos embeddings for 100 residues modulo $2$ through $101$,
        capturing periodicity and number-theoretic structure.
\end{itemize}
The two streams are fused via FiLM (Feature-wise Linear Modulation)~\cite{perez2018film}.
During training, three predictors are jointly optimised in a multi-task objective: magnitude
regression, sign classification, and modulo prediction for 100 moduli. Concrete integer values
are recovered from the masked-position predictions (magnitude, sign, residue distributions)
using a probabilistic Chinese Remainder Theorem (CRT)-based \textbf{Solver} (Section~\ref{sec:solver}).

This combination of the modulo-spectrum representation and the CRT solver is a
\emph{neural-symbolic} design: residue arithmetic and the Chinese Remainder Theorem are
built into the model and the decoding procedure as explicit number-theoretic structure,
rather than left for the network to discover end-to-end. Our contribution is therefore not
the prediction accuracy itself---predicting integer sequences via machine learning is a
long-standing task, e.g.\ the Kaggle Integer Sequence Learning
competition~\cite{kaggle-integer-seq}---but the representation that explicitly encodes
number-theoretic structure, together with quantitative evidence that this representation
captures the structure of OEIS sequences. At the same time, our results document how
difficult direct neural prediction of OEIS terms remains even with this structure built in,
providing motivation for hybrid symbolic--neural approaches (Section~\ref{sec:conclusion}).

We evaluate IntSeqBERT against two baselines on 274,705 OEIS sequences across three model
sizes (Small: 6.4M, Middle: 29.0M, Large: 91.5M parameters), all on a
single GeForce RTX 3070 Ti (8\,GB VRAM).
The implementation, preprocessing scripts, and evaluation code are publicly
available.\footnote{\url{https://github.com/aabaa/IntSeqBERT}}

\noindent\textbf{Contributions.}
\begin{enumerate}
  \item \textit{IntSeqBERT architecture}: a dual-stream Transformer fusing magnitude and
        modular-arithmetic embeddings via FiLM, jointly trained on OEIS with magnitude
        regression, sign classification, and 100-dimensional modulo prediction. At the
        Large scale, IntSeqBERT achieves \textbf{95.85\%} magnitude accuracy and
        \textbf{50.38\%} MMA, surpassing Vanilla by $+8.9$\,pt and $+4.5$\,pt, respectively.
        The dual-stream representation yields a \textbf{7.4-fold} improvement in
        Solver-based next-term prediction (Top-1: 19.09\% vs.\ 2.59\%).

  \item \textit{Number-theoretic finding}: modulo spectrum analysis reveals that NIG is
        strongly negatively correlated with $\varphi(m)/m$ (where $\varphi$ denotes Euler's totient function; Pearson correlation of $r = {-}0.851$ ($p < 10^{-28}$)),
        providing quantitative evidence that composite moduli aggregate arithmetic structure
        more efficiently via CRT.

  \item \textit{Scaling behaviour}: modulo accuracy and Solver accuracy improve more steeply
        with model size than magnitude accuracy, suggesting arithmetic reasoning benefits
        disproportionately from increased capacity.
\end{enumerate}

\noindent\textbf{Paper organisation.}
Sections~\ref{sec:background}--\ref{sec:setup} review related work, present the architecture,
and describe the experimental setup. Section~\ref{sec:experiments} reports results, Section~\ref{sec:analysis} provides analyses, and Section~\ref{sec:conclusion} concludes.
